% installation.tex
%
% The documentation in this file is part of Pyxplot
% <http://www.pyxplot.org.uk>
%
% Copyright (C) 2006-2026 Dominic Ford <coders@pyxplot.org.uk>
%               2008-2013 Ross Church
%
% $Id$
%
% Pyxplot is free software; you can redistribute it and/or modify it under the
% terms of the GNU General Public License as published by the Free Software
% Foundation; either version 2 of the License, or (at your option) any later
% version.
%
% You should have received a copy of the GNU General Public License along with
% Pyxplot; if not, write to the Free Software Foundation, Inc., 51 Franklin
% Street, Fifth Floor, Boston, MA  02110-1301, USA

% ----------------------------------------------------------------------------

% LaTeX source for the Pyxplot Users' Guide

\chapter{Installation}

\label{ch:installation}

In this chapter we describe how to install Pyxplot on a range of UNIX-like
operating systems.

Pyxplot works on most UNIX-like operating systems. We have tested it under
Linux and MacOS\index{MacOS}, and believe that it
should work on other similar POSIX systems. We regret that it is not available
for Microsoft Windows, and have no plans for porting it at this time.

\section{Installation within Linux distributions}

Debian\index{Debian Linux}\index{installation!under Debian}\footnote{See
\url{https://packages.debian.org/pyxplot}} and Ubuntu\index{Ubuntu
Linux}\index{installation!under Ubuntu}\footnote{See
\url{https://packages.ubuntu.com/pyxplot}} include Pyxplot in their package
archives, though at the time of writing they offer Pyxplot~0.9.2. Please note
that this manual describes Pyxplot~0.10.x. To install the packaged version of
Pyxplot under Debian or Ubuntu, simply type:

\begin{verbatim}
apt-get install pyxplot gv
\end{verbatim}

Users of other distributions, or who want a newer version of Pyxplot, should
use the {\tt .tar.gz} archives available from the Pyxplot website. The process
is described below.

\section{System requirements}

Pyxplot depends on the following software packages, which are not included in
the source tarball:\index{system requirements}

\vspace{0.5cm}
\begin{itemize}
\item fftw (version 2 or, preferably, 3+) \index{fftw}
\item gcc and make\index{gcc}\index{make}
\item Ghostscript \index{Ghostscript}
\item The Gnu Scientific Library (version 1.10+) \index{gsl}
\item ImageMagick \index{ImageMagick}
\item \latexdcf\ (version $2\epsilon$; a full installation is likely to be required in distributions which offer a choice) \index{latex}
\item libpng (version 1.2+) \index{libpng}
\item libxml2 (version 2.6+) \index{libxml}
\item pkg-config (or pkgconf, which provides it) \index{pkg-config}
\item zlib \index{zlib}
\end{itemize}
\vspace{0.5cm}

\noindent It is very strongly recommended that the following software packages
also be installed:

\vspace{0.5cm}
\begin{itemize}
\item cfitsio -- required for Pyxplot to be able to plot \datafile s in FITS format.
\item Ghostview \index{Ghostview} ({\tt gv}) -- required for Pyxplot to be able to display plots live on the screen; Pyxplot remains able to generate image files on disk without it. Alternatively, the \indcmdt{set viewer} within Pyxplot allows a different PostScript viewer to be used.
\item gunzip \index{gunzip} -- required for Pyxplot to be able to plot compressed \datafile s in {\tt .gz} format.
\item The Gnu Readline Library (version 5+) \index{readline} -- required for Pyxplot to be able to provide tab completion and command histories in Pyxplot's interactive command-line interface.
\item libkpathsea \index{libkpathsea} -- required to efficiently find the fonts used by \latexdcf. Without it, Pyxplot uses the {\tt kpsewhich} program, which is part of TeX Live, instead.
\item wget \index{wget} -- required for Pyxplot to be able to plot \datafile s directly from the Internet.
\end{itemize}
\vspace{0.5cm}

\noindent In the case of the recommended packages, Pyxplot tests for the
availability of each when it is installed and issues a warning if any are not
found. Installation can proceed, but some of Pyxplot's features will be
disabled. If they are later added to the system, Pyxplot should be reinstalled
to take advantage of their presence.

If the C compiler supports OpenMP\index{OpenMP}, as {\tt gcc} does, Pyxplot
uses it to sample functions on several processor cores at once.

\subsection{Dependencies in Debian and Ubuntu}

Debian and Ubuntu users can find the above software in the following
packages\footnote{The package names listed here were checked against Ubuntu
26.04, and are also used by recent Debian releases. However, packages
occasionally change name between versions.}: \index{Debian Linux}\index{Ubuntu
Linux}\index{installation!under Debian}\index{installation!under Ubuntu}

\vspace{2mm}
\noindent {\tt gcc}, {\tt ghostscript}, {\tt gv}, {\tt imagemagick}, {\tt libc6-dev}, {\tt libcfitsio-dev},\newline
\noindent {\tt libfftw3-dev}, {\tt libgsl-dev}, {\tt libkpathsea-dev}, {\tt libpng-dev},\newline
\noindent {\tt libreadline-dev}, {\tt libxml2-dev}, {\tt make}, {\tt pkgconf}, {\tt texlive-fonts-extra},\newline
\noindent {\tt texlive-fonts-recommended}, {\tt texlive-latex-extra},\newline
\noindent {\tt texlive-latex-recommended}, {\tt wget}, {\tt zlib1g-dev}.
\vspace{2mm}

\noindent These packages may be installed from a command prompt by typing, all on one line:

\begin{verbatim}
sudo apt-get install gcc ghostscript gv imagemagick libc6-dev
       libcfitsio-dev libfftw3-dev libgsl-dev libkpathsea-dev
       libpng-dev libreadline-dev libxml2-dev make pkgconf
       texlive-fonts-extra texlive-fonts-recommended
       texlive-latex-extra texlive-latex-recommended wget
       zlib1g-dev
\end{verbatim}

\subsection{Displaying plots on the screen in Ubuntu 26.04}
\index{Ubuntu Linux}\index{AppArmor}

Users of Ubuntu 26.04 need to change the system's AppArmor security settings
for Pyxplot to be able to display plots on the screen. Ubuntu 26.04 confines
Ghostscript with an AppArmor profile ({\tt /etc/apparmor.d/gs}) which does not
allow it to connect to the X display. \ghostview\ relies on Ghostscript to draw
its pages, so Pyxplot's {\tt X11} terminals show only a blank page; output to
files, e.g.\ in {\tt eps}, {\tt pdf} or {\tt png} format, is not affected. The
kernel log reports the problem with messages such as:

\begin{verbatim}
apparmor="DENIED" operation="connect" profile="gs" ...
          peer_addr="@/tmp/.X11-unix/X0"
\end{verbatim}

\noindent To allow Ghostscript to connect to the X display, while keeping the
profile's other restrictions, add the following line to the local additions to
the profile, and then reload it, by typing:

\begin{verbatim}
echo 'include <abstractions/X>' | sudo tee -a /etc/apparmor.d/local/gs
sudo apparmor_parser -r /etc/apparmor.d/gs
\end{verbatim}

\noindent Alternatively, Pyxplot can be told to use a viewer which does not run
the Ghostscript executable, e.g.\ by typing {\tt set viewer okular} within
Pyxplot (see the {\tt set viewer} command in Chapter~\ref{ch:reference}).

\subsection{Dependencies in Fedora}
\index{Fedora Linux}\index{installation!under Fedora}

Fedora users can install the above software by typing, all on one
line\footnote{The package names listed here were checked against Fedora~43
and~44.}:

\begin{verbatim}
sudo dnf install cfitsio-devel fftw-devel gcc ghostscript
       gsl-devel gv ImageMagick libpng-devel libxml2-devel make
       pkgconf-pkg-config readline-devel texlive-latex
       texlive-lib-devel wget2-wget zlib-ng-compat-devel
\end{verbatim}

\noindent If \latexdcf\ reports that packages or fonts are missing, a larger
TeX Live scheme, e.g.\ {\tt texlive-scheme-medium}, should be installed.

\subsection{Dependencies in MacOS}
\index{MacOS X}\index{MacPorts}

Users of MacOS can find the above software in the following MacPorts packages:

\vspace{2mm}
\noindent {\tt cfitsio}, {\tt fftw-3}, {\tt ghostscript}, {\tt gsl}, {\tt gv}, {\tt ImageMagick}, {\tt libpng},\newline
\noindent {\tt libxml2}, {\tt pkgconfig}, {\tt readline}, {\tt texlive}, {\tt wget}, {\tt zlib}.
\vspace{2mm}

It may then be necessary to run the command
\begin{verbatim}
export C_INCLUDE_PATH=/opt/local/include LIBRARY_PATH=/opt/local/lib
\end{verbatim}
before running the {\tt configure} script below. \ghostview\ is an X11
application, and needs an X server such as XQuartz to display plots.

Alternatively, users of Homebrew\index{Homebrew} can install the software by
typing:

\begin{verbatim}
brew install cfitsio fftw ghostscript gsl imagemagick libpng
       libxml2 pkgconf readline texlive wget
\end{verbatim}

\noindent and then, before running the {\tt configure} script below:

\begin{verbatim}
export C_INCLUDE_PATH=$(brew --prefix)/include
export LIBRARY_PATH=$(brew --prefix)/lib
\end{verbatim}

\noindent Homebrew does not package \ghostview, so Pyxplot cannot display plots
on the screen unless another PostScript viewer is selected with the
\indcmdt{set viewer}; it can still write plots to files. Apple's C compiler does
not support OpenMP, so Pyxplot samples functions on a single processor core.

\section{Installation from source archive}
\index{installation}

First, download the required archive from the front page of the
Pyxplot website -- \url{http://www.pyxplot.org.uk}. It is assumed that the
packages listed above have already been installed; if they are not, you will
need to either install them yourself, if you have superuser access to your
machine, or contact your system administrator.\index{installation!user-level}

\begin{itemize}
\item Unpack the distributed .tar.gz:

\begin{verbatim}
tar xvfz pyxplot_0.10.0.tar.gz
cd pyxplot-0.10.0
\end{verbatim}

\item Run the installation script:

\begin{verbatim}
./configure
make
\end{verbatim}

\item Finally, start Pyxplot:

\begin{verbatim}
./bin/pyxplot
\end{verbatim}

\end{itemize}

\subsection{System-wide installation}

Having completed the steps described above, Pyxplot may be installed
system-wide by a superuser with the following additional
step:\index{installation!system-wide}

\begin{verbatim}
sudo make install
\end{verbatim}

By default, the Pyxplot executable installs to {\tt /usr/local/bin/pyxplot}.
If desired, this installation path may be modified in the file {\tt
Makefile.skel}, by changing the variable {\tt USRDIR} near the top of the file to an
alternative desired installation location.

Pyxplot may now be started by any user of the system, simply by typing:

\begin{verbatim}
pyxplot
\end{verbatim}

